单部位实验中，仅胸部的测试 Macro-F1 最高，为 0.7690。完整三部位随机森林为 0.9261；移除实验中移除手腕的测试 F1 最低，为 0.8132。完整配置比最佳单部位高 15.71 个百分点。三组部位对照均使用相同的训练、验证和测试窗口，输入维度分别为 50、100 或 150。

以完整三部位为参照，移除手腕、胸部和脚踝分别降低 11.29、1.00 和 4.95 个百分点。手腕输入在这一特征组合中带来的成绩增量最大。完整配置也高于各双部位配置，说明三个部位的统计表示在本次识别任务中提供了互补信息。